\chapter{Fungsi Holomorfik}\label{ch:b3:holomorphic}

Keterturunan kompleks tampak seperti ragam kecil atas
teori realnya --- satu limit, satu hasil bagi. Padahal ia
semesta yang berbeda. Sebuah fungsi yang dapat diturunkan \emph{sekali} pada
himpunan bagian terbuka $\C$ dengan sendirinya dapat diturunkan tak berhingga kali,
bersifat analitik, tertentukan pada seluruh daerahnya oleh nilainya di dekat
satu titik, dan terkendala oleh asas global yang kaku
(Liouville, modulus maksimum). Semua ini mengalir dari satu
keajaiban, yakni teorema Cauchy: bahwa integral sebuah \omterm{def:b3:holomorphic:holo}{fungsi
holomorfik} sepanjang \omterm{def:b3:topology:pathconnected}{lintasan} tertutup di daerah berbentuk bintang lenyap.
Bab ini membuktikan keajaiban itu (lewat argumen Goursat, tanpa
keteraturan yang diandaikan selain keterturunannya), memanen
akibat klasiknya, lalu menuntaskan sebuah teorema yang buku ini
pakai secara utang sejak \cref{ch:b3:galois}: bahwa setiap
polinomial kompleks yang tak konstan berakar.

Sepanjang bab ini, $\Omega \subseteq \C$ terbuka, dan $D(a, r)$
menandai cakram terbukanya.

\section{Keterturunan kompleks}

\begin{definition}\label{def:b3:holomorphic:holo}
Sebuah $f \colon \Omega \to \C$ disebut \emph{holomorfik}\index{fungsi
holomorfik} pada $\Omega$ jika untuk setiap $z_0 \in \Omega$
\[
f'(z_0) = \lim_{h\to0}\frac{f(z_0 + h) - f(z_0)}{h}
\]
ada ($h \in \C^*$). Jumlah, hasil kali, hasil bagi (dengan penyebut
yang tak lenyap), dan komposisi fungsi holomorfik bersifat
holomorfik, dengan rumus yang lazim (sebab bukti Tahun ke-1 dan ke-2 berlaku
kata demi kata: karena ia hanya memakai operasi lapangan dan limit).
Sedangkan $\mathcal H(\Omega)$ menandai himpunan fungsi holomorfik
pada $\Omega$.
\end{definition}

\begin{proposition}[Cauchy--Riemann]
\label{prop:b3:holomorphic:cauchyriemann}
Tulis $f(x + \iu y) = P(x,y) + \iu Q(x,y)$. Maka $f$ bersifat
\omterm{def:b3:holomorphic:holo}{holomorfik} di $z_0$ jika dan hanya jika $f$ dapat diturunkan-$\R$ di $z_0$
(sebagai pemetaan dua variabel real) dan\index{persamaan
Cauchy--Riemann}
\[
\frac{\partial P}{\partial x} = \frac{\partial Q}{\partial y},
\qquad
\frac{\partial P}{\partial y} = -\frac{\partial Q}{\partial x}
\qquad \text{di } z_0 ;
\]
setara dengan itu, diferensial realnya adalah perkalian dengan
bilangan kompleks $f'(z_0)$.
\end{proposition}

\begin{proof}
Keterturunan-$\C$-nya mengatakan $f(z_0 + h) = f(z_0) + ch +
o(\abs h)$ dengan $c = f'(z_0)$: yakni sebuah diferensial linear-$\R$
yang berupa perkalian dengan $c = a + \iu b$, yakni bermatriks
$\bigl(\begin{smallmatrix} a & -b\\ b & a
\end{smallmatrix}\bigr)$ pada basis $(1, \iu)$ --- tepat
hubungan yang ditampilkan bagi turunan parsialnya. Sebaliknya diferensial
semacam itu bersifat linear-$\C$, dan definisi $o(\abs h)$-nya
bersesuaian.
\end{proof}

\begin{example}\label{ex:b3:holomorphic:examples}
Polinomial dalam $z$, fungsi rasional di luar kutubnya, dan
--- menurut teorema penurunan suku demi suku Tahun ke-2 untuk deret
pangkat, yang buktinya bekerja secara identik atas $\C$ --- setiap jumlah
deret pangkat $\sum a_n(z - a)^n$ di dalam cakram
kekonvergenannya: semuanya \omterm{def:b3:holomorphic:holo}{holomorfik}, dengan turunan $\sum na_n(z -
a)^{n-1}$ (berjari-jari sama). Khususnya $\exp z = \sum z^n/n!$
bersifat utuh (yakni \omterm{def:b3:holomorphic:holo}{holomorfik} pada $\C$) dengan $\exp' = \exp$. Sedangkan
di sisi lain $z \mapsto \bar z$, $\abs z$, dan $\operatorname{Re}z$
tak \omterm{def:b3:holomorphic:holo}{holomorfik} di mana pun (sebab Cauchy--Riemann gagal di mana-mana):
jadi \omterm{def:b3:holomorphic:holo}{keholomorfikan} adalah kekakuan yang mengawetkan arah dan sudut, bukan
kemulusan.
\end{example}

\section{Integral kontur}

\begin{definition}\label{def:b3:holomorphic:contour}
Sebuah \emph{\omterm{def:b3:topology:pathconnected}{lintasan}} adalah pemetaan $\mathcal C^1$ sepenggal-sepenggal $\gamma \colon
\intcc ab \to \C$; dan ia \emph{tertutup} jika $\gamma(a) =
\gamma(b)$. Untuk $f$ yang \omterm{def:b3:topology:continuity}{kontinu} pada peta
$\gamma$:\index{integral kontur}
\[
\int_\gamma f(z)\,\dd z =
\int_a^b f(\gamma(t))\,\gamma'(t)\,\dd t,
\qquad
\Bigl|\int_\gamma f\,\dd z\Bigr| \leq
\sup_{\gamma}\abs f\cdot\operatorname{length}(\gamma)
\]
(yakni \emph{ketaksamaan ML}; dengan panjangnya $= \int_a^b\abs{\gamma'}$).
Integralnya invarian terhadap reparametrisasi $\mathcal C^1$ yang
naik dan berganti tanda saat arahnya dibalik.
\end{definition}

\begin{proposition}[Antiturunan]\label{prop:b3:holomorphic:primitive}
Untuk $f$ yang \omterm{def:b3:topology:continuity}{kontinu} pada $\Omega$, pernyataan berikut setara:
(i) $f$ mempunyai \emph{antiturunan} $F \in \mathcal H(\Omega)$ (yakni $F'
= f$); (ii) $\int_\gamma f\,\dd z = 0$ untuk setiap \omterm{def:b3:topology:pathconnected}{lintasan} tertutup
$\gamma$ di $\Omega$. Dalam hal itu $\int_\gamma f\,\dd z =
F(\gamma(b)) - F(\gamma(a))$ untuk setiap \omterm{def:b3:topology:pathconnected}{lintasan}.
\end{proposition}

\begin{proof}
(i)$\Rightarrow$: $\frac{\dd}{\dd t}F(\gamma(t)) =
F'(\gamma(t))\gamma'(t)$ (lewat aturan rantai, yang sah sepenggal-sepenggal), sehingga
integralnya berteleskop menjadi selisih ujungnya; dan \omterm{def:b3:topology:pathconnected}{lintasan} tertutup
memberikan $0$. (ii)$\Rightarrow$(i): tetapkan $z_*$ pada sebuah \omterm{def:b3:topology:components}{komponen
terhubung}, lalu definisikan $F(z) = \int_{\gamma_z}f\,\dd z$ sepanjang sembarang
\omterm{def:b3:topology:pathconnected}{lintasan} dari $z_*$ ke $z$ (yang terdefinisi baik: sebab dua \omterm{def:b3:topology:pathconnected}{lintasan} berselisih
sebuah \omterm{def:b3:topology:pathconnected}{lintasan} tertutup); lalu untuk $h$ yang kecil, dengan mengambil ruas dari $z$ ke
$z + h$,
\[
\frac{F(z + h) - F(z)}{h} - f(z)
= \frac1h\int_{[z, z+h]}\bigl(f(w) - f(z)\bigr)\dd w
\longrightarrow 0
\]
lewat ketaksamaan ML dan \omterm{def:b3:topology:continuity}{kekontinuan} $f$ di $z$.
\end{proof}

\begin{definition}[Bilangan lilitan]\label{def:b3:holomorphic:index}
Untuk sebuah \omterm{def:b3:topology:pathconnected}{lintasan} tertutup $\gamma$ dan $z \notin
\operatorname{im}\gamma$, \emph{indeks}\index{bilangan
lilitan}\index{indeks sebuah lintasan} itu adalah
\[
\operatorname{Ind}_\gamma(z) = \frac1{2\iu\pi}
\int_\gamma\frac{\dd w}{w - z} .
\]
Ia sebuah \emph{bilangan bulat}: sebab dengan menetapkan $\varphi(t) =
\int_a^t\frac{\gamma'(s)}{\gamma(s) - z}\dd s$, fungsi
$(\gamma(t) - z)\eu^{-\varphi(t)}$ berturunan nol
(sepenggal-sepenggal), sehingga konstan; lalu di $t = b$, $\eu^{\varphi(b)} =
\frac{\gamma(b) - z}{\gamma(a) - z} = 1$, sehingga $\varphi(b) \in
2\iu\pi\Z$. Sebagai fungsi $z$, indeksnya \omterm{def:b3:topology:continuity}{kontinu} pada
$\C\setminus\operatorname{im}\gamma$ (lewat kekonvergenan terdominasi),
sehingga konstan pada setiap \omterm{def:b3:topology:components}{komponen terhubungnya}, dan $0$ pada
komponen yang tak terbatas (lewat ML: sebab integralnya menuju $0$ saat $z \to
\infty$). Untuk lingkaran $\gamma(t) = a + r\eu^{\iu t}$ dengan $t
\in \intcc0{2\pi}$: $\operatorname{Ind}_\gamma(z) = 1$ bagi $z
\in D(a,r)$ (hitunglah di $z = a$: $\frac1{2\iu\pi}
\int_0^{2\pi}\frac{r\iu\eu^{\iu t}}{r\eu^{\iu t}}\dd t = 1$;
lalu ketetapannya menyelesaikan sisanya).
\end{definition}

\section{Teorema Cauchy}

\begin{theorem}[Goursat]\label{thm:b3:holomorphic:goursat}
Misalkan $f \in \mathcal H(\Omega)$ dan $T \subseteq \Omega$ sebuah
segitiga padat tertutup. Maka $\int_{\partial T}f\,\dd z = 0$
(dengan batasnya dilalui sekali, dalam arah mana pun).\index{teorema Goursat}
\end{theorem}

\begin{proof}
Misalkan $I(T) = \int_{\partial T}f\,\dd z$. Menghubungkan titik tengah
sisinya memecah $T$ menjadi empat segitiga separuh \omterm{def:b3:measure:measure}{ukuran} $T^{(1)},
\dots, T^{(4)}$, dan rusuk dalamnya saling meniadakan berpasangan: sehingga $I(T) =
\sum_iI(T^{(i)})$. Pilihlah $T_1$ di antaranya yang
$\abs{I(T_1)} \geq \frac14\abs{I(T)}$, lalu iterasikan: diperoleh barisan
bersarang $T \supseteq T_1 \supseteq T_2 \supseteq\cdots$ dengan
\[
\abs{I(T_n)} \geq 4^{-n}\abs{I(T)},
\qquad
\operatorname{diam}T_n = 2^{-n}\operatorname{diam}T,
\quad
\operatorname{length}(\partial T_n) =
2^{-n}\operatorname{length}(\partial T).
\]
Irisannya $\bigcap T_n$ berupa satu titik $z_0$
(sebab \omterm{def:b3:topology:compact}{kompak} bersarang yang diameternya melenyap,
\cref{thm:b3:topology:compactprops}(3)). Untuk keterturunannya di
$z_0$: diberikan $\varepsilon$, maka untuk $n$ yang besar, pada $T_n$,
\[
f(z) = f(z_0) + f'(z_0)(z - z_0) + R(z),
\qquad \abs{R(z)} \leq \varepsilon\abs{z - z_0} \leq
\varepsilon\operatorname{diam}T_n .
\]
Bagian afinnya mempunyai antiturunan: sehingga integralnya pada
$\partial T_n$ yang tertutup lenyap (\cref{prop:b3:holomorphic:primitive}),
menyisakan
\[
\abs{I(T_n)} = \Bigl|\int_{\partial T_n}R\Bigr|
\leq \varepsilon\operatorname{diam}(T_n)\,
\operatorname{length}(\partial T_n)
= \varepsilon\,4^{-n}\operatorname{diam}(T)
\operatorname{length}(\partial T) .
\]
Dengan membandingkannya terhadap $\abs{I(T_n)} \geq 4^{-n}\abs{I(T)}$:
$\abs{I(T)} \leq \varepsilon\cdot\text{konstanta}$ untuk setiap
$\varepsilon$: sehingga $I(T) = 0$.
\end{proof}

\begin{theorem}[Teorema Cauchy, versi berbentuk bintang]
\label{thm:b3:holomorphic:cauchy}
Misalkan $\Omega$ \emph{berbentuk bintang} terhadap $c$ (yakni setiap ruas
$[c, z]$ dengan $z \in \Omega$ terletak di $\Omega$) --- misalnya cembung.
Maka setiap $f \in \mathcal H(\Omega)$ mempunyai antiturunan pada $\Omega$;
sehingga $\int_\gamma f\,\dd z = 0$ untuk \emph{setiap}
\omterm{def:b3:topology:pathconnected}{lintasan} tertutup $\gamma$ di $\Omega$.\index{teorema Cauchy}
\end{theorem}

\begin{proof}
Definisikan $F(z) = \int_{[c,z]}f\,\dd w$. Untuk $z, z + h \in
\Omega$ dengan $[z, z+h] \subseteq \Omega$ (yang benar untuk $h$ yang kecil),
segitiga bertitik sudut $c, z, z+h$ terletak di $\Omega$
(lewat kebintangannya: sebab setiap titiknya terletak pada sebuah ruas $[c,
w]$ dengan $w \in [z, z+h] \subseteq \Omega$): sehingga Goursat memberikan
\[
F(z + h) - F(z) = \int_{[z, z+h]}f\,\dd w,
\]
dan perhitungan hasil bagi selisih pada
\cref{prop:b3:holomorphic:primitive} melahirkan $F' = f$. Sedangkan
kelenyapan integral \omterm{def:b3:topology:pathconnected}{lintasan} tertutupnya menyusul dari proposisi
yang sama.
\end{proof}

\begin{theorem}[Rumus integral Cauchy]
\label{thm:b3:holomorphic:formula}
Misalkan $f \in \mathcal H(\Omega)$, $\bar D(a, r) \subseteq
\Omega$, dan $C_r$ lingkaran $\partial D(a,r)$ yang dilalui sekali
berlawanan arah jarum jam. Maka untuk setiap $z \in D(a, r)$:\index{rumus
integral Cauchy}
\[
f(z) = \frac1{2\iu\pi}\int_{C_r}\frac{f(w)}{w - z}\,\dd w .
\]
\end{theorem}

\begin{proof}
Tetapkan $z$ lalu definisikan pada $\Omega$
\[
g(w) = \begin{cases}
\dfrac{f(w) - f(z)}{w - z} & w \neq z,\\[4pt]
f'(z) & w = z :
\end{cases}
\]
maka $g$ \omterm{def:b3:topology:continuity}{kontinu} pada $\Omega$ dan \omterm{def:b3:holomorphic:holo}{holomorfik} di luar $z$. Lalu Goursat
berlaku bagi $g$ pada setiap segitiga $T \subseteq \Omega'$, dengan
$\Omega'$ sebuah cakram yang sedikit lebih besar daripada $\bar D(a,r)$ di dalam
$\Omega$: sebab jika $z \notin T$, secara langsung; sedangkan jika $z \in T$, pecahlah $T$
menjadi segitiga kecil bertitik sudut $z$ ditambah segitiga
yang menghindari $z$; lalu pada segitiga bertitik sudut $z$ batas ML memberikan
$\abs{\int_{\partial T'}g} \leq \sup_{T'}\abs
g\cdot\operatorname{length} \to 0$ saat segitiganya menyusut,
sedangkan kepingan sisanya lenyap menurut Goursat --- sehingga
$\int_{\partial T}g = 0$ dalam segala kasusnya. Sedangkan bukti
\cref{thm:b3:holomorphic:cauchy} hanya memakai sifat segitiga ini:
sehingga $g$ mempunyai antiturunan pada $\Omega'$ yang cembung, jadi
$\int_{C_r}g = 0$, yakni
\[
\frac1{2\iu\pi}\int_{C_r}\frac{f(w)}{w - z}\dd w
= f(z)\,\frac{1}{2\iu\pi}\int_{C_r}\frac{\dd w}{w - z}
= f(z)\operatorname{Ind}_{C_r}(z) = f(z) .
\]
\end{proof}

\section{Keanalitikan dan runtunannya}

\begin{theorem}[Holomorfik $=$ analitik]
\label{thm:b3:holomorphic:analytic}
Misalkan $f \in \mathcal H(\Omega)$ dan $D(a, R) \subseteq \Omega$.
Maka\index{keanalitikan fungsi holomorfik}
\[
f(z) = \sum_{n\geq0}c_n\,(z - a)^n \quad \text{pada } D(a, R),
\qquad
c_n = \frac1{2\iu\pi}\int_{C_r}\frac{f(w)}{(w -
a)^{n+1}}\,\dd w \ \ (0 < r < R),
\]
dengan koefisiennya tak bergantung pada $r$. Karena itu $f$ dapat
diturunkan-$\C$ tak berhingga kali, $c_n = f^{(n)}(a)/n!$, dan
\emph{taksiran Cauchy}\index{taksiran Cauchy} berlaku:
\[
\abs{c_n} \;\leq\; \frac{\sup_{\abs{w - a} =
r}\abs{f(w)}}{r^{n}} .
\]
\end{theorem}

\begin{proof}
Untuk $\abs{z - a} < r$: uraikan kernel Cauchynya menjadi
deret geometri
\[
\frac1{w - z} = \frac1{(w - a)\bigl(1 - \frac{z - a}{w -
a}\bigr)} = \sum_{n\geq0}\frac{(z - a)^n}{(w - a)^{n+1}},
\]
yang konvergen normal terhadap $w$ pada $C_r$ (sebab $\abs{\frac{z-a}{w-a}} =
\frac{\abs{z-a}}r < 1$): lalu integralkan suku demi suku terhadap
$\frac{f(w)}{2\iu\pi}$ (sebab kekonvergenan seragamnya membenarkan
penukarannya) lalu terapkan \cref{thm:b3:holomorphic:formula}. Sedangkan
deret pangkat dapat diturunkan tak berhingga kali dengan $c_n =
f^{(n)}(a)/n!$ (Tahun ke-2), yang juga menunjukkan bahwa $c_n$ tak
bergantung pada $r$. Untuk taksirannya: batasi integral koefisiennya
lewat ML.
\end{proof}

\begin{example}[Kesingularannya menentukan jari-jarinya]
\label{ex:b3:holomorphic:radius}
Mengapa fungsi real yang tak berbahaya $\frac1{1 + x^2}$ mempunyai
deret Taylor di $x = 3$ yang konvergen hanya untuk $\abs{x - 3} <
\sqrt{10}$, padahal tak ada yang salah pada garis realnya? Sebab
teorema di atas membuat jari-jari kekonvergenan di $a$
sama dengan jarak dari $a$ ke titik terdekat yang
\omterm{def:b3:holomorphic:holo}{keholomorfikannya} gagal. Di sini $f(z) = \frac1{1 + z^2}$ bersifat
\omterm{def:b3:holomorphic:holo}{holomorfik} tepat pada $\C\setminus\{\pm\iu\}$, sehingga
uraiannya di $a = 3$ konvergen pada cakram terbesar yang menghindari
$\pm\iu$, yakni berjari-jari $\abs{3 - \iu} = \sqrt{10}$ --- dan
ia tak dapat konvergen pada cakram yang lebih besar, sebab jumlahnya akan memperluas
$f$ secara \omterm{def:b3:holomorphic:holo}{holomorfik} ke sebuah \omterm{def:b3:topology:topology}{persekitaran} $\pm\iu$, tempat
$\abs f \to \infty$. Jadi teori realnya melihat jari-jari misterius
$\sqrt{10}$; sedangkan bidang kompleksnya melihat dua kutub.
Inilah kaidah praktisnya: \emph{untuk mencari jari-jari
kekonvergenan, temukanlah kesingularannya} --- misalnya deret
Taylor $\tan$ di $0$ berjari-jari $\frac\pi2$ (yakni akar
$\cos$ yang terdekat), dan fungsi pembangkit Bernoulli
$\frac z{\eu^z - 1}$ (\cref{pb:b3:holomorphic:1}, Bagian VI)
berjari-jari $2\pi$ (yakni akar tak nol $\eu^z - 1$ yang terdekat:
$\pm2\iu\pi$).
\end{example}

\begin{corollary}[Liouville; d'Alembert--Gauss]
\label{cor:b3:holomorphic:liouville}
Sebuah fungsi utuh yang terbatas bersifat konstan. Karena itu setiap
polinomial tak konstan atas $\C$ berakar: yakni \emph{$\C$ bersifat
\omterm{def:b3:galois:closure}{tertutup secara aljabar}}.\index{teorema Liouville}\index{teorema
fundamental aljabar}
\end{corollary}

\begin{proof}
Jika $\abs f \leq M$ pada $\C$: maka untuk setiap $a$ dan $r$, $\abs{c_1(a)}
= \abs{f'(a)} \leq M/r \to 0$: sehingga $f' \equiv 0$, dan $f$
konstan (pada $\C$ yang \omterm{def:b3:topology:connected}{terhubung}: sebab turunan nol mengakibatkan
kekonstanan lokal --- lalu integralkan sepanjang ruasnya). Jika $P$ tak
berakar, maka $1/P$ akan utuh dan terbatas (sebab $\abs{P(z)} \to
\infty$ saat $\abs z \to \infty$: karena suku utamanya mendominasi, sehingga
$\abs{1/P}$ kecil di luar sebuah cakram besar dan \omterm{def:b3:topology:continuity}{kontinu} pada
cakram \omterm{def:b3:topology:compact}{kompaknya}): jadi konstan --- yang mustahil bagi $P$ yang tak konstan.
(Sedangkan soal akhir pekannya memberikan bukti kedua yang dasar beserta
akibat aljabarnya.)
\end{proof}

\begin{theorem}[Akarnya terpencil; teorema keidentikan]
\label{thm:b3:holomorphic:identity}
Misalkan $\Omega$ \emph{\omterm{def:b3:topology:connected}{terhubung}} dan $f \in \mathcal
H(\Omega)$ dengan $f \not\equiv 0$. Maka setiap akar $a$ dari $f$ ber\emph{orde}
berhingga: yakni $f(z) = (z - a)^m\,g(z)$ dengan $g \in
\mathcal H(\Omega)$ dan $g(a) \neq 0$, dan akar $f$ tak
bertitik tumpuk di $\Omega$. Karena itu, jika dua
\omterm{def:b3:holomorphic:holo}{fungsi holomorfik} pada $\Omega$ bersesuaian pada sebuah himpunan yang bertitik
tumpuk di $\Omega$, keduanya bersesuaian
di mana-mana.\index{teorema keidentikan}\index{akar terpencil}
\end{theorem}

\begin{proof}
Misalkan $Z$ himpunan titik yang semua turunan $f$-nya
lenyap. Maka $Z$ tertutup (sebab irisan himpunan tertutup) dan terbuka:
sebab jika semua $c_n = 0$ di $a$, maka uraian deret pangkatnya membuat $f
\equiv 0$ pada sebuah cakram di sekitar $a$. Lewat keterhubungannya: $Z =
\varnothing$ atau $Z = \Omega$; sedangkan yang terakhir dikecualikan oleh $f
\not\equiv 0$. Jadi pada sebuah akar $a$, suatu koefisiennya tak nol:
misalkan $m$ minimal dengan $c_m \neq 0$; maka $f(z) = (z -
a)^m\sum_{n\geq m}c_n(z-a)^{n-m}$ pada sebuah cakram, dan jumlahnya
mendefinisikan $g$ yang \omterm{def:b3:holomorphic:holo}{holomorfik} di dekat $a$ dengan $g(a) = c_m \neq 0$;
lalu perluas $g = f/(z-a)^m$ di luar $a$ (yang \omterm{def:b3:holomorphic:holo}{holomorfik} di sana). Karena
$g(a) \neq 0$ dan $g$ \omterm{def:b3:topology:continuity}{kontinu}, maka $f$ tak berakar lain pada
sebuah \omterm{def:b3:topology:topology}{persekitaran} $a$: sehingga akarnya terpencil, dan himpunan berisi titik
terpencil tak bertitik tumpuk di $\Omega$ (sebab sebuah
titik tumpuk akarnya adalah akar --- lewat \omterm{def:b3:topology:continuity}{kekontinuannya} --- dan
tak akan terpencil). Untuk keidentikannya: terapkanlah pada selisihnya,
yang himpunan akarnya bertitik tumpuk, sehingga memaksanya ke
cabang $Z = \Omega$.
\end{proof}

\begin{theorem}[Nilai rata-rata dan modulus maksimum]
\label{thm:b3:holomorphic:maximum}
Misalkan $f \in \mathcal H(\Omega)$.
\begin{enumerate}
  \item (Nilai rata-rata) Untuk $\bar D(a, r) \subseteq \Omega$:
        $f(a) = \frac1{2\pi}\int_0^{2\pi}f(a +
        r\eu^{\iu t})\,\dd t$.
  \item (Asas maksimum) Jika $\Omega$ \omterm{def:b3:topology:connected}{terhubung} dan
        $\abs f$ mencapai maksimum lokal di suatu titik
        $\Omega$, maka $f$ konstan. Karena itu, untuk
        $\Omega$ yang terbatas dan $f$ yang \omterm{def:b3:topology:continuity}{kontinu} pada
        $\bar\Omega$: $\sup_{\bar\Omega}\abs f =
        \sup_{\partial\Omega}\abs f$.\index{asas maksimum}
\end{enumerate}
\end{theorem}

\begin{proof}
(1) adalah rumus Cauchy di pusatnya: parametrikanlah $C_r$.
(2) Katakanlah $\abs f \leq \abs{f(a)}$ pada $\bar D(a, \rho)
\subseteq \Omega$. Jika $f(a) = 0$, maka $f \equiv 0$ di dekat $a$.
Selain itu, untuk $0 < r \leq \rho$, nilai rata-ratanya memberikan
\[
\abs{f(a)} \leq \frac1{2\pi}\int_0^{2\pi}\abs{f(a +
r\eu^{\iu t})}\,\dd t \leq \abs{f(a)} :
\]
sehingga integran \omterm{def:b3:topology:continuity}{kontinu} tak negatifnya $\abs{f(a)} - \abs{f(a +
r\eu^{\iu t})}$ berrata-rata nol, jadi lenyap: sehingga $\abs f$ bernilai
\emph{konstan} $= \abs{f(a)} \ne 0$ pada cakramnya. Sedangkan \omterm{def:b3:holomorphic:holo}{fungsi
holomorfik} yang bermodulus tetap tak nol pada sebuah cakram bersifat konstan:
sebab menurunkan $P^2 + Q^2 = \text{konstanta}$ memberikan $PP_x + QQ_x
= 0$ dan $PP_y + QQ_y = 0$; lalu menyubstitusikan hubungan Cauchy--Riemann
$P_y = -Q_x$ dan $Q_y = P_x$ ke persamaan keduanya
melahirkan sistem linear
\[
P\,P_x + Q\,Q_x = 0, \qquad -P\,Q_x + Q\,P_x = 0,
\]
yang determinannya $P^2 + Q^2 \neq 0$: sehingga $P_x = Q_x = 0$, jadi
$f' = P_x + \iu Q_x = 0$ pada cakramnya: sehingga $f$ konstan di sana. Lalu
teorema keidentikannya menyebarkan kekonstanan itu ke seluruh $\Omega$. Untuk bentuk batasnya: $\abs f$
mencapai supremumnya pada $\bar\Omega$ yang \omterm{def:b3:topology:compact}{kompak}; sedangkan maksimum
di \omterm{def:b3:topology:interior}{interiornya} membuat $f$ konstan, sehingga supremumnya tercapai pada
batasnya dalam segala kasus.
\end{proof}

\begin{theorem}[Teorema kekonvergenan Weierstrass]
\label{thm:b3:holomorphic:weierstrassconv}
Jika $f_n \in \mathcal H(\Omega)$ konvergen ke $f$ secara seragam pada
setiap himpunan bagian \omterm{def:b3:topology:compact}{kompak} $\Omega$, maka $f \in \mathcal
H(\Omega)$ dan $f_n^{(k)} \to f^{(k)}$ secara seragam pada \omterm{def:b3:topology:compact}{kompak},
untuk setiap $k$.\index{teorema kekonvergenan Weierstrass}
\end{theorem}

\begin{proof}
Fungsi $f$ \omterm{def:b3:topology:continuity}{kontinu}. Untuk sembarang segitiga tertutup $T \subseteq
\Omega$: $\int_{\partial T}f = \lim\int_{\partial T}f_n = 0$
(lewat kekonvergenan seragam pada $\partial T$ yang \omterm{def:b3:topology:compact}{kompak}; dan Goursat untuk
$f_n$). Lalu menurut argumen \cref{thm:b3:holomorphic:cauchy},
$f$ mempunyai antiturunan lokal (sebab cakramnya cembung; dan hanya sifat
segitiganya yang dipakai), yakni $f = F'$ secara lokal dengan $F$ yang
\omterm{def:b3:holomorphic:holo}{holomorfik}; sedangkan $F$ bersifat analitik
(\cref{thm:b3:holomorphic:analytic}), sehingga demikian pula $f = F'$:
jadi \omterm{def:b3:holomorphic:holo}{holomorfik}. (Inilah \emph{teorema
Morera}\index{teorema Morera}: bahwa \omterm{def:b3:topology:continuity}{kontinu} dengan integral
segitiga yang lenyap mengakibatkan \omterm{def:b3:holomorphic:holo}{holomorfik}.) Untuk turunannya: bagi
$\bar D(a, 2r) \subseteq \Omega$ dan $z \in \bar D(a, r)$,
rumus Cauchy untuk turunannya (turunkanlah
\cref{thm:b3:holomorphic:formula} di bawah integralnya, atau pakailah
rumus koefisiennya) memberikan
\[
\abs{f_n'(z) - f'(z)} =
\Bigl|\frac{1}{2\iu\pi}\int_{C_{2r}}\frac{f_n(w) - f(w)}{(w -
z)^2}\,\dd w\Bigr|
\leq \frac{2r\,\sup_{C_{2r}}\abs{f_n - f}}{r^2} \to 0
\]
secara seragam pada $\bar D(a,r)$; lalu selimuti sebuah \omterm{def:b3:topology:compact}{kompak} oleh berhingga banyak
cakram semacam itu, dan iterasikan untuk $k$ yang lebih tinggi.
\end{proof}

\begin{omfigure}
\begin{tikzpicture}[scale=1.5]
  \draw[->] (-1.7,0) -- (1.9,0);
  \draw[->] (0,-1.35) -- (0,1.5);
  \draw[omDef, very thick] (0,0) circle (1);
  \foreach \a in {0, 30, 60, 90, 120, 150}
    \draw[omDef!50] (\a:1.06) -- (\a:0.94)
      (\a+180:1.06) -- (\a+180:0.94);
  \fill[omThm] (0.35,0.25) circle (1pt)
    node[above, font=\small] {$z$};
  \fill (0,0) circle (1pt) node[below left, font=\small] {$a$};
  \draw[omProp, thick, ->]
    (1,0) arc (0:75:1);
  \node[omProp, font=\small] at (60:1.28) {$C_r$};
  \node[font=\small, omThm] at (1.45,0.75)
    {$f(z) = \dfrac1{2\iu\pi}\displaystyle
     \int_{C_r}\dfrac{f(w)}{w-z}\,\dd w$};
\end{tikzpicture}

{\small Rumus Cauchy: bahwa nilai sebuah \omterm{def:b3:holomorphic:holo}{fungsi holomorfik}
\emph{di dalam} sebuah cakram merupakan rata-rata berbobot nilainya pada
lingkaran pembatasnya. Segala yang kaku tentang \omterm{def:b3:holomorphic:holo}{keholomorfikan} ---
keanalitikan, Liouville, asas maksimum --- terbentang dari
satu identitas ini.}
\end{omfigure}

\begin{method}\label{met:b3:holomorphic:toolkit}
Perkakas hariannya. Untuk membuktikan sebuah \omterm{def:b3:holomorphic:holo}{fungsi holomorfik}: tampilkanlah ia
sebagai deret pangkat, sebuah komposisi, sebuah limit seragam lokal
(\cref{thm:b3:holomorphic:weierstrassconv}), atau sebuah \omterm{thm:b3:lebesgue:paramcont}{integral
berparameter} \omterm{def:b3:holomorphic:holo}{holomorfik} (\cref{exo:b3:holomorphic:7} ---
turunkanlah di bawah $\int$ atau terapkan Morera--Fubini). Untuk membuktikan
identitas: buktikanlah pada sebuah ruas atau subdaerah lalu panggillah
teorema keidentikannya. Untuk membatasi: pakailah \omterm{thm:b3:holomorphic:analytic}{taksiran Cauchy} pada
lingkaran terbesar yang tersedia. Untuk membuktikan kekonstanan atau ketakadaan:
pakailah Liouville atau asas maksimumnya. Ketahuilah selalu \emph{di mana}
fungsi Anda \omterm{def:b3:holomorphic:holo}{holomorfik} dan \emph{cakram mana} yang muat di
$\Omega$.
\end{method}

\section{Latihan}

\begin{exercise}[$\star$]\label{exo:b3:holomorphic:1}
(a) Di titik mana saja $z \mapsto \bar z$, $\abs z^2$, dan
$\operatorname{Re}z$ dapat diturunkan secara kompleks? \omterm{def:b3:holomorphic:holo}{Holomorfik} pada sebuah
\omterm{def:b3:topology:topology}{himpunan terbuka}?
(b) Tunjukkan bahwa $P(x, y) = x^2 - y^2$ merupakan bagian real sebuah
\omterm{def:b3:holomorphic:holo}{fungsi holomorfik} pada $\C$, yang dicari secara eksplisit, lalu tentukanlah
\emph{semuanya}.
\end{exercise}

\begin{exercise}[$\star$]\label{exo:b3:holomorphic:2}
Hitunglah dari definisinya: $\int_{C}z^n\,\dd z$ untuk setiap $n
\in \Z$ dengan $C$ lingkaran satuannya; lalu $\int_\gamma\bar z\,\dd z$ sepanjang
ruas $[0, 1+\iu]$ dan sepanjang \omterm{def:b3:topology:pathconnected}{lintasan} dua ruas
yang melalui $1$: lalu simpulkan bahwa $\bar z$ tak berantiturunan pada
\omterm{def:b3:topology:topology}{persekitaran} mana pun atas \omterm{def:b3:topology:pathconnected}{lintasan} ini.
\end{exercise}

\begin{exercise}[$\star\star$]\label{exo:b3:holomorphic:3}
(a) Tunjukkan bahwa logaritma pokoknya $\log z = \ln\abs z +
\iu\arg z$ (dengan $\arg \in \intoo{-\pi}\pi$) bersifat \omterm{def:b3:holomorphic:holo}{holomorfik} pada
$\C\setminus\intoc{-\infty}0$ dengan turunan $\frac1z$
\emph{(yakni antiturunan $\frac1z$ pada bidang teriris yang berbentuk bintang:
\cref{thm:b3:holomorphic:cauchy}; lalu tetapkan konstantanya)}.
(b) Tunjukkan bahwa tak ada logaritma \omterm{def:b3:topology:continuity}{kontinu} pada
$\C^*$ \emph{(dengan rintangan bebas turunannya: yakni \omterm{def:b3:holomorphic:index}{indeks}
lingkaran satuannya)}.
(c) Uraikan $\log(1 + z)$ menjadi deret pangkat pada $D(0,1)$.
\end{exercise}

\begin{exercise}[$\star\star$]\label{exo:b3:holomorphic:4}
(a) Misalkan $f$ utuh dengan $\abs{f(z)} \leq C(1 + \abs
z)^{n}$. Tunjukkan bahwa $f$ merupakan polinomial berderajat $\leq n$
\emph{(lewat \omterm{thm:b3:holomorphic:analytic}{taksiran Cauchy} pada lingkaran yang besar)}.
(b) Misalkan $f$ utuh dengan $\operatorname{Re}f$ yang terbatas dari atas.
Tunjukkan bahwa $f$ konstan \emph{(tinjaulah $\eu^{f}$)}.
(c) Turunkan ``Picard kecil untuk pemetaan afin'': bahwa fungsi
utuh yang melewatkan sebuah setengah bidang bersifat konstan.
\end{exercise}

\begin{exercise}[$\star\star$]\label{exo:b3:holomorphic:5}
(a) Misalkan $f$ \omterm{def:b3:holomorphic:holo}{holomorfik} pada $\Omega \ni 0$ yang \omterm{def:b3:topology:connected}{terhubung} dengan
$f(\frac1n) = \frac1{n^2}$ untuk setiap $n$ yang besar. Tentukanlah $f$.
(b) Adakah suatu $f$ \omterm{def:b3:holomorphic:holo}{holomorfik} pada $\C^*$ yang memenuhi $f(\frac1n) =
\frac{(-1)^n}{n}$ untuk setiap $n \geq 1$? Benarkanlah.
(c) Tampilkan dua \omterm{def:b3:holomorphic:holo}{fungsi holomorfik} berbeda pada
$D(0,1)\sqcup D(3,1)$ yang bersesuaian pada $D(0,1)$: lalu di manakah
teorema keidentikannya memakai keterhubungan?
\end{exercise}

\begin{exercise}[$\star\star$]\label{exo:b3:holomorphic:6}
Misalkan $f$ \omterm{def:b3:holomorphic:holo}{holomorfik} pada cakram satuan terbuka $\mathbb D$,
\omterm{def:b3:topology:continuity}{kontinu} pada $\bar{\mathbb D}$, dengan $\abs f \equiv 1$ pada
lingkaran batasnya.
(a) Jika $f$ tak berakar di $\mathbb D$, tunjukkan bahwa $f$ konstan
\emph{(terapkan asas maksimumnya pada $f$ dan pada $1/f$)}.
(b) Berikan sebuah contoh berakar yang $f$-nya tak konstan.
\end{exercise}

\begin{exercise}[$\star\star$]\label{exo:b3:holomorphic:7}
(\omterm{def:b3:holomorphic:holo}{Keholomorfikan} di bawah integralnya) Misalkan $\mu$ \omterm{def:b3:measure:measure}{ukuran} berhingga
pada sebuah ruang $X$ dan $g \colon X\times\Omega \to \C$ dengan:
$g(x, \cdot) \in \mathcal H(\Omega)$ untuk setiap $x$, $g$
\omterm{def:b3:lebesgue:measurable}{terukur} terhadap $x$, dan $\abs g \leq h(x)$ dengan $h$ yang \omterm{def:b3:lebesgue:l1}{terintegralkan},
secara seragam lokal terhadap $z$. Tunjukkan bahwa $G(z) = \int_Xg(x, z)\dd\mu(x)$
bersifat \omterm{def:b3:holomorphic:holo}{holomorfik} pada $\Omega$. \emph{(Lewat Morera: sebab integral segitiganya
lenyap menurut Fubini dan Goursat; sedangkan \omterm{def:b3:topology:continuity}{kekontinuannya} lewat kekonvergenan
terdominasi. Lalu terapkan pada $\Gamma(z) =
\int_0^\infty t^{z-1}\eu^{-t}\dd t$ pada $\{\operatorname{Re}z >
0\}$.)}
\end{exercise}

\begin{exercise}[$\star\star\star$]\label{exo:b3:holomorphic:8}
(Gauss--Lucas) Misalkan $P \in \C[X]$ tak konstan. Tunjukkan bahwa
setiap akar $P'$ terletak di selubung cembung akar
$P$. \emph{(Tulislah $\frac{P'}{P} = \sum_k\frac{m_k}{z - a_k}$
di sebuah akar $z$ dari $P'$ yang bukan akar $P$, lalu ambil
konjugatnya, lalu bacalah kombinasi cembungnya.)} Peragakan pada
$P = z^3 - 1$.
\end{exercise}

\begin{exercise}[$\star\star\star$]\label{exo:b3:holomorphic:9}
Misalkan $f$ utuh dan berperiode ganda: $f(z + 1) = f(z + \iu)
= f(z)$ untuk setiap $z$. Tunjukkan bahwa $f$ konstan. \emph{(Batasilah
$f$ pada bujur sangkar fundamentalnya yang \omterm{def:b3:topology:compact}{kompak}, lalu di mana-mana;
lalu Liouville.)} Moralnya: bahwa fungsi eliptik yang tak konstan haruslah
berkutub --- yakni tema \cref{ch:b3:residues}.
\end{exercise}

\begin{exercise}[$\star\star$]\label{exo:b3:holomorphic:10}
(a) Tunjukkan bahwa $P = \operatorname{Re}f$ dari sebuah $f$ yang \omterm{def:b3:holomorphic:holo}{holomorfik}
memenuhi sifat nilai rata-rata $P(a) =
\frac1{2\pi}\int_0^{2\pi}P(a + r\eu^{\iu t})\dd t$ dan bersifat
\emph{harmonik}: yakni $\partial^2_{xx}P + \partial^2_{yy}P = 0$
(turunkanlah Cauchy--Riemann; lalu pakailah
\cref{thm:b3:holomorphic:analytic} untuk kemulusan yang diperlukan).
(b) Turunkan asas maksimumnya bagi bagian real \omterm{def:b3:holomorphic:holo}{fungsi
holomorfik} pada daerah yang terbatas.
\end{exercise}

\begin{exercise}[$\star\star\star$]\label{exo:b3:holomorphic:11}
(Pencerminan Schwarz) Misalkan $\Omega^+ = \{z : \abs z < 1,\
\operatorname{Im}z > 0\}$, $I = \intoo{-1}1$, dan $f$
\omterm{def:b3:holomorphic:holo}{holomorfik} pada $\Omega^+$, \omterm{def:b3:topology:continuity}{kontinu} pada $\Omega^+\cup I$,
serta \emph{bernilai real pada $I$}. Definisikan
\[
F(z) = \begin{cases} f(z) & z \in \Omega^+\cup I,\\
\overline{f(\bar z)} & \bar z \in \Omega^+ . \end{cases}
\]
(a) Tunjukkan bahwa $F$ terdefinisi baik dan \omterm{def:b3:topology:continuity}{kontinu} pada
$\Omega = \Omega^+\cup I\cup\Omega^-$, serta \omterm{def:b3:holomorphic:holo}{holomorfik} pada
$\Omega^\pm$ \emph{(untuk $\Omega^-$: periksalah Cauchy--Riemann
bagi $\overline{f(\bar z)}$, atau uraikan $f$ menjadi deret pangkat
lokal lalu konjugatkan koefisiennya)}.
(b) Tunjukkan bahwa $F$ \omterm{def:b3:holomorphic:holo}{holomorfik} pada seluruh $\Omega$ lewat
kriteria Morera: yakni $\int_{\partial T}F = 0$ untuk setiap
segitiga $T \subseteq \Omega$ \emph{(pecahlah segitiganya di $I$
lalu dorong sisi mendatarnya menjauh dari sumbunya sejauh
$\varepsilon$, dengan memakai \omterm{def:b3:topology:continuity}{kekontinuan} seragamnya)}.
(c) Turunkan: bahwa \omterm{def:b3:holomorphic:holo}{fungsi holomorfik} pada cakram yang real pada sebuah
diameter memenuhi $f(\bar z) = \overline{f(z)}$; dan bahwa
\omterm{def:b3:holomorphic:holo}{fungsi holomorfik} yang tak konstan tak dapat bernilai real pada
himpunan bagian terbuka tak kosong mana pun dari daerahnya (yang \omterm{def:b3:topology:connected}{terhubung}).
\end{exercise}

\begin{exercise}[$\star\star$]\label{exo:b3:holomorphic:12}
(Persamaan Pythagoras kompleks) Carilah semua pasangan fungsi
\emph{utuh} dengan $f^2 + g^2 = 1$.
(a) Tunjukkan bahwa $h = f + \iu g$ bersifat utuh dan bebas akar, dan
bahwa setiap fungsi utuh yang bebas akar berbentuk $\eu^{\varphi}$ untuk
suatu $\varphi$ yang utuh \emph{(sebab $h'/h$ utuh, sehingga berantiturunan
pada $\C$ yang berbentuk bintang; lalu sesuaikan konstantanya lalu
tunjukkan bahwa $h\eu^{-\varphi}$ konstan)}.
(b) Simpulkan $f = \cos\varphi$ dan $g = \sin\varphi$ dengan
$\varphi$ yang utuh, lalu periksalah konversnya. Apakah
solusi utuh $f^2 + g^2 = 0$?
\end{exercise}

\section{Soal: teorema fundamental aljabar, dua kali}

\begin{problem}[{Soal akhir pekan --- $\C$ tertutup secara
aljabar: bukti d'Alembert, bukti Liouville, dan
panennya}]\label{pb:b3:holomorphic:1}
Misalkan $P(z) = z^n + a_{n-1}z^{n-1} + \dots + a_0$ dengan $n \geq 1$.
Kita membuktikan dua kali bahwa $P$ berakar, lalu memungut apa yang telah
ditunggu aljabar sejak \cref{ch:b3:galois}.

\medskip
\noindent\textbf{Bagian I --- Kekoersifan dan minimumnya.}
\begin{enumerate}
  \item Tunjukkan bahwa $\abs{P(z)} \to +\infty$ saat $\abs z \to
        \infty$: tepatnya, $\abs{P(z)} \geq \frac12\abs z^n$
        untuk $\abs z \geq R_0$ yang sesuai.
  \item Turunkan bahwa $\abs P$ mencapai minimum global pada $\C$:
        yakni ada $z_0$ dengan $\abs{P(z_0)} = \inf_\C\abs P$
        \emph{(lewat kekompakan sebuah cakram tertutup yang besar,
        \cref{cor:b3:topology:heineborel})}.
\end{enumerate}

\medskip
\noindent\textbf{Bagian II --- Penurunan d'Alembert.} Andaikan,
demi pertentangannya, $P(z_0) \neq 0$.
\begin{enumerate}[resume]
  \item Uraikan $Q(h) = P(z_0 + h)/P(z_0)$ sebagai polinomial dalam
        $h$: $Q(h) = 1 + c_kh^k + h^{k+1}S(h)$ dengan $c_k \neq
        0$, $k \geq 1$, dan $S$ sebuah polinomial.
  \item Pilihlah arah penurunannya: untuk $t > 0$ yang kecil,
        tetapkan $h = t\,\omega$ dengan $\omega^k = -1/c_k$ (mengapa
        $\omega$ semacam itu ada? --- \emph{buktikanlah}
        keberadaan akar ke-$k$ sembarang bilangan kompleks lewat
        bentuk kutubnya, secara bebas dari teorema yang sedang
        dibuktikan). Lalu tunjukkan bahwa
        \[
        \abs{Q(t\omega)} \leq 1 - t^k + C\,t^{k+1}
        \]
        untuk $t$ yang kecil, dengan konstanta $C$ yang eksplisit.
  \item Simpulkan $\abs{Q(t\omega)} < 1$ untuk $t$ yang kecil ---
        yang bertentangan dengan keminimalan $\abs{P(z_0)}$. Karena itu
        $P(z_0) = 0$: yakni \emph{setiap polinomial kompleks
        yang tak konstan berakar} (d'Alembert--Argand).
\end{enumerate}

\medskip
\noindent\textbf{Bagian III --- Satu baris Liouville,
selengkapnya.}
\begin{enumerate}[resume]
  \item Tuliskan dengan saksama bukti
        \cref{cor:b3:holomorphic:liouville}: bahwa jika $P$ tak
        berakar, periksalah bahwa $1/P$ utuh dan terbatas (ukurlah,
        dengan memakai pertanyaan 1), sehingga konstan, lalu simpulkan.
        Lalu bandingkan kedua buktinya: bahan mana yang dipakai masing-masing?
        (Kekompakannya muncul pada keduanya --- di mana?)
\end{enumerate}

\medskip
\noindent\textbf{Bagian IV --- Panennya.}
\begin{enumerate}[resume]
  \item Tunjukkan bahwa setiap $P \in \C[X]$ berderajat $n$ terbelah:
        $P = c\prod_{i}(X - \alpha_i)^{m_i}$ dengan $\sum m_i =
        n$ \emph{(lewat induksi dan pembagian Euclid oleh $(X -
        \alpha)$)}.
  \item Tunjukkan bahwa polinomial \omterm{def:b3:rings:divisibility}{tak tereduksi} $\R[X]$ adalah
        yang linear dan yang kuadratik berdiskriminan negatif
        \emph{(dengan memasangkan akar sekawannya)}; lalu turunkan
        bahwa setiap polinomial real berderajat gasal berakar real,
        lalu berikan bukti kedua yang berdasar urutan bagi
        fakta terakhir itu (lewat teorema nilai antara) ---
        sambil memeriksa bahwa keduanya bersesuaian pada $X^3 - X - 1$.
  \item Turunkan utang yang kini dapat dibayar buku ini: (i) bahwa setiap
        endomorfisma ruang vektor-$\C$ berdimensi berhingga
        yang tak nol bernilai eigen, sehingga setiap matriks
        kompleks mempunyai \omterm{thm:b3:modules:jordan}{bentuk Jordan}
        (\cref{thm:b3:modules:jordan}); (ii) bahwa lapangan
        $\bar\Q$ berisi bilangan aljabar yang dipakai pada
        \cref{rem:b3:galois:closureexamples} sungguh merupakan \omterm{def:b3:galois:closure}{penutup
        aljabar} $\Q$.
  \item (\omterm{def:b3:topology:interior}{Penutup}) Tunjukkan dengan tepat di mana setiap buktinya akan patah atas
        lapangan seperti $\Q(\iu)$: langkah mana yang memakai keberadaan
        akar ke-$k$ (pertanyaan 4), dan mana yang memakai kekompakan
        atau \omterm{def:b3:complete:complete}{kelengkapan} (pertanyaan 2 dan 6)? Lalu simpulkan dalam lima
        baris: bahwa teoremanya sungguh bersifat \emph{analitik} ---
        sebab setiap buktinya di suatu tempat memanggil \omterm{def:b3:complete:complete}{kelengkapan} atau
        keterhubungan $\R$ --- walaupun pernyataannya
        murni aljabar.
\end{enumerate}

\medskip
\noindent\textbf{Bagian V --- Tangga kekakuan fungsi
utuh.} Liouville adalah anak tangga pertamanya; dan kita
menaikinya.
\begin{enumerate}[resume]
  \item (\omterm{thm:b3:holomorphic:analytic}{Taksiran Cauchy}) Dari rumus Cauchy pada
        lingkaran berjari-jari $r$ di sekitar $a$, buktikan
        \[
        \bigl|f^{(n)}(a)\bigr| \leq
        \frac{n!\,\sup_{\abs{z-a}=r}\abs f}{r^n} ,
        \]
        lalu pulihkan Liouville sebagai kasus $n = 1$ dengan $r \to
        \infty$.
  \item (Pertumbuhan polinomial) Tunjukkan bahwa $f$ yang utuh dengan
        $\abs{f(z)} \leq A + B\abs z^m$ untuk setiap $z$ merupakan
        polinomial berderajat $\leq m$ \emph{(bunuhlah koefisien
        Taylornya di luar $m$ dengan pertanyaan 11)}.
  \item (Bagian real yang terbatas) Tunjukkan bahwa $f$ yang utuh dengan
        $\operatorname{Re}f$ terbatas dari atas bersifat konstan
        \emph{(terapkan Liouville pada $\eu^{f}$)}.
  \item (Keperiodikan ganda) Misalkan $f$ utuh dengan $f(z + 1)
        = f(z)$ dan $f(z + \iu) = f(z)$ untuk setiap $z$. Tunjukkan bahwa $f$
        konstan. Lalu simpulkan: bahwa fungsi ``eliptik'' yang tak konstan
        haruslah berkesingularan --- yakni alasan historis
        masuknya kutub ke analisis kompleks.
  \item (Jangkauan yang padat) Tunjukkan bahwa jangkauan sebuah fungsi utuh
        yang tak konstan bersifat padat di $\C$: sebab jika $f(\C)$ melewatkan sebuah
        cakram $D(a, r)$, maka $\frac1{f - a}$ utuh dan
        terbatas. (Picard membuktikan bahwa jangkauannya melewatkan paling banyak
        \emph{satu titik}; sedangkan \omterm{ex:b3:lebesgue:gamma}{kepadatan} adalah aras yang dicapai perkakas kita.)
  \item (Meledak di tak hingga $\Rightarrow$ polinomial) Andaikan $f$
        utuh dan $\abs{f(z)} \to \infty$ saat $\abs z \to
        \infty$. Tunjukkan: bahwa akar $f$ berhingga banyaknya
        ($z_1, \dots, z_p$, dengan kelipatan $m_i$); bahwa
        hasil baginya $g = f/\prod(z - z_i)^{m_i}$ utuh dan
        bebas akar; bahwa $1/g$ bertumbuh polinomial, sehingga
        (menurut pertanyaan 12) merupakan polinomial, yang tentu konstan
        (sebab bebas akar); lalu simpulkan bahwa $f$ merupakan polinomial. Jadi
        di antara fungsi utuh, polinomiallah tepat yang
        meledak di tak hingga --- sedangkan $\eu^z$ gagal meledak sepanjang
        $\R_-$.
\end{enumerate}

\medskip
\noindent\textbf{Bagian VI --- Bayangan harmonik dan sebuah nilai
rata-rata Gauss.}
\begin{enumerate}[resume]
  \item Misalkan $f = u + \iu v$ \omterm{def:b3:holomorphic:holo}{holomorfik} pada sebuah \omterm{def:b3:topology:topology}{himpunan terbuka}.
        Periksalah bahwa $u = \operatorname{Re}f$ memenuhi
        sifat nilai rata-rata
        \[
        u(a) = \frac1{2\pi}\int_0^{2\pi}
        u\bigl(a + r\eu^{\iu\theta}\bigr)\,\dd\theta
        \]
        (yakni bagian real rumus Cauchynya), lalu turunkan
        asas maksimumnya bagi $u$ pada daerah yang terbatas, dengan
        bukti keterhubungan yang sama seperti bagi $\abs f$.
  \item (Nilai rata-rata Gauss) Untuk $a \in \C$ dan $r > 0$ dengan
        $\abs a \neq r$, buktikan bahwa
        \[
        \frac1{2\pi}\int_0^{2\pi}
        \log\bigl|a - r\eu^{\iu\theta}\bigr|\,\dd\theta
        = \log\max\bigl(\abs a, r\bigr)
        \]
        \emph{(jika $\abs a > r$: maka $z \mapsto \log\abs{a - z}$
        merupakan bagian real sebuah logaritma \omterm{def:b3:holomorphic:holo}{holomorfik} pada sebuah
        \omterm{def:b3:topology:topology}{persekitaran} cakram tertutupnya --- mengapa ada yang demikian? ---
        sehingga pertanyaan 17 berlaku; sedangkan jika $\abs a < r$:
        faktorkanlah $\abs{a - r\eu^{\iu\theta}} = r\,\abs{1 -
        \frac ar\eu^{-\iu\theta}}$ lalu pakailah kembali kasus
        pertamanya)}.
  \item (\omterm{def:b3:measure:measure}{Ukuran} Mahler) Untuk $P = c\prod_{i=1}^n(X -
        \alpha_i) \in \C[X]$, turunkan \emph{rumus Jensen
        bagi polinomial}:
        \[
        \frac1{2\pi}\int_0^{2\pi}\log\bigl|P(\eu^{\iu\theta})
        \bigr|\,\dd\theta
        = \log\Bigl(\abs c\prod_{i}\max(1,
        \abs{\alpha_i})\Bigr) :
        \]
        yakni rata-rata ukur $\abs P$ pada lingkaran satuannya
        membaca akarnya yang di luar cakram. Lalu periksalah pada $P =
        X^2 - X$ dan pada $P = 2X - 1$.
  \item (Bilangan Bernoulli) Definisikan koefisien $B_n$ lewat
        $\frac{z}{\eu^z - 1} = \sum_{n\geq0}\frac{B_n}{n!}
        z^n$ di dekat $0$ (mengapa ruas kirinya analitik di
        $0$?). Turunkan rekurensinya
        $\sum_{k=0}^{n}\binom{n+1}kB_k = 0$ (untuk $n \geq 1$) dari
        $(\eu^z - 1)\cdot\frac z{\eu^z-1} = z$, lalu hitunglah $B_0,
        \dots, B_6$, dan tunjukkan $B_{2k+1} = 0$ untuk $k \geq 1$
        \emph{(sebab fungsi $\frac z{\eu^z-1} + \frac z2$ bersifat
        genap)}. Bilangan ini akan menghargai setiap $\zeta(2k)$ di
        \cref{ch:b3:residues}.
  \item (Kerealan) Tunjukkan bahwa fungsi utuh yang bernilai real
        pada $\R$ memenuhi $f(\bar z) =
        \overline{f(z)}$ di mana-mana \emph{(bandingkan koefisien
        Taylornya di $0$, atau terapkan teorema keidentikannya
        pada $z \mapsto \overline{f(\bar z)}$)}; lalu turunkan lagi
        bahwa akar tak real dari polinomial real muncul berpasangan
        sekawan (yakni pemasangan pertanyaan 8, yang dibuktikan kembali
        secara analitik).
  \item (Moralnya) Rakitlah tangga kekakuannya: terbatas
        $\Rightarrow$ konstan; terbatas secara polinomial
        $\Rightarrow$ polinomial; meledak di tak hingga $\Rightarrow$
        polinomial; melewatkan sebuah cakram $\Rightarrow$ konstan;
        berperiode ganda $\Rightarrow$ konstan. Lalu bandingkan dalam
        sebuah paragraf pendek dengan fungsi $\mathcal C^\infty$ real
        (yakni \omterm{def:b3:lp:mollifier}{fungsi gundukan},
        \cref{thm:b3:lp:regularization}): mengapa \omterm{def:b3:holomorphic:holo}{keholomorfikan},
        yang merupakan syarat murni lokal, memberlakukan hukum dan
        ketertiban yang global?
\end{enumerate}

\medskip
\noindent\textbf{Bagian VII --- Panen terakhir.}
\begin{enumerate}[resume]
  \item (Ketaksamaan Landau) Untuk $P = \sum_{k=0}^na_kX^k$,
        buktikan nilai rata-ratanya $\frac1{2\pi}\int_0^{2\pi}
        \abs{P(\eu^{\iu\theta})}^2\dd\theta =
        \sum_k\abs{a_k}^2$ (lewat keortogonalan
        $\eu^{\iu k\theta}$), lalu, dengan memakai batas titik demi titiknya
        $\log t \leq t - 1$ untuk membandingkan rata-rata
        $\log\abs P^2$ dan $\abs P^2$, turunkan dari pertanyaan
        19 bahwa
        \[
        \abs c\prod_{i}\max\bigl(1, \abs{\alpha_i}\bigr)
        \;\leq\;
        \Bigl(\sum_{k=0}^{n}\abs{a_k}^2\Bigr)^{1/2} :
        \]
        yakni hasil kali akar di luar cakram satuannya
        terkendali oleh \omterm{def:b3:measure:measure}{ukuran} $\ell^2$ koefisiennya
        (tanganilah akar yang \emph{pada} lingkarannya dengan menerapkan
        ketaksamaannya pada $P(rX)$ lalu membiarkan $r \to 1$). Lalu periksalah
        pada $X^2 - X$.
  \item (Bilangan Bernoulli tumbuh secara faktorial) Tunjukkan bahwa
        jari-jari kekonvergenan $\sum\frac{B_n}{n!}z^n$ bernilai
        \emph{tepat} $2\pi$: setidaknya $2\pi$ sebab
        $z/(\eu^z - 1)$ diperluas secara \omterm{def:b3:holomorphic:holo}{holomorfik} ke
        $D(0, 2\pi)$, dan paling banyak $2\pi$ sebab jumlahnya jika tidak
        akan tetap terbatas di dekat $2\pi\iu$, tempat
        $\abs{z/(\eu^z-1)} \to \infty$. Lalu turunkan
        \[
        \limsup_{k\to\infty}
        \Bigl(\frac{\abs{B_{2k}}}{(2k)!}\Bigr)^{1/2k}
        = \frac1{2\pi} :
        \]
        yakni bilangan Bernoulli tumbuh secara faktorial. Lalu dengan menerima $B_{12}
        = -\frac{691}{2730}$ (dari rekurensi pertanyaan 20,
        yang dijalankan lebih jauh), bandingkan $\abs{B_{12}} \approx 0.25311$
        dengan ramalan yang lebih tajam $2\,(2k)!/(2\pi)^{2k}
        \approx 0.25305$ pada $k = 6$ --- yakni empat angka yang cocok
        dari sebuah hukum asimtotik yang akan dibuktikan
        \cref{ch:b3:residues} secara persis, lewat $\zeta(2k)$.
  \item (Akarnya bergerak secara \omterm{def:b3:topology:continuity}{kontinu}) Misalkan $(P_j)$ monik
        berderajat $n$ dengan koefisien yang konvergen ke koefisien
        $P$ yang (monik). Buktikan \emph{batas Cauchy}: bahwa setiap
        akar sebuah $Q = X^n + \sum_{k<n}q_kX^k$ yang monik
        memenuhi $\abs\alpha \leq 1 + \max_k\abs{q_k}$;
        lalu turunkan bahwa akar $P_j$ tetap berada di sebuah himpunan
        \omterm{def:b3:topology:compact}{kompak} yang tetap, dan, dengan memetik subbarisan
        konvergen vektor akarnya lalu mengambil
        limitnya pada pemfaktoran pertanyaan 7, bahwa
        multihimpunan akar $P_j$ konvergen ke multihimpunan akar $P$.
        Lalu tunjukkan akhirnya bahwa \omterm{def:b3:topology:continuity}{kekontinuan} adalah yang terbaik yang dapat
        dikatakan: sebab untuk $P_\varepsilon = X^2 - 2X + 1 +
        \varepsilon$, gangguan berukuran $\varepsilon =
        10^{-4}$ menggeser akar rangkap $1$ sejauh $10^{-2}$ ---
        yakni eksponen Hölder $\frac1m$ pada akar berlipat $m$,
        dan tak pernah Lipschitz: sehingga secara numerik, akar berlipat berbiaya
        separuh angkanya.

\end{enumerate}
\end{problem}
